
\documentclass{article}
\usepackage{colortbl}
\usepackage{makecell}
\usepackage{multirow}
\usepackage{supertabular}

\begin{document}

\newcounter{utterance}

\centering \large Interaction Transcript for game `bbh', experiment `formal\_fallacies', episode 158 with reflexion{-}llama{-}3.1{-}8b{-}instruct.
\vspace{24pt}

{ \footnotesize  \setcounter{utterance}{1}
\setlength{\tabcolsep}{0pt}
\begin{supertabular}{c@{$\;$}|p{.15\linewidth}@{}p{.15\linewidth}p{.15\linewidth}p{.15\linewidth}p{.15\linewidth}p{.15\linewidth}}
    \# & \multicolumn{2}{c}{Player} && \multicolumn{2}{c}{Game Master} \\
    \hline

    \theutterance \stepcounter{utterance}  
    & & & \multicolumn{4}{p{0.6\linewidth}}{
        \cellcolor[rgb]{0.9,0.9,0.9}{
            \makecell[{{p{\linewidth}}}]{
                \texttt{\tiny{[P1$\langle$GM]}}
                \texttt{You have attempted this same task before and did not succeed. These are the plans you wrote after those attempts.} \\
\texttt{Trial \#1: New plan: } \\
\\ 
\texttt{I failed to identify the validity of the fourth argument. Upon re{-}examining the argument, I realized that the conclusion "Being an ancestor of Janice is necessary for being a close friend of Francine" does not logically follow from the premises. The premises state that every close friend of Francine is a cousin of Samantha or an ancestor of Janice or a close friend of Shanika, and that being a close friend of Shanika is necessary for being a close friend of Francine. However, this does not imply that being an ancestor of Janice is necessary for being a close friend of Francine. In fact, the premises do not provide any information about the relationship between being a close friend of Francine and being an ancestor of Janice. Therefore, the argument is deductively invalid. Next time, I will carefully analyze the premises and the conclusion to determine the validity of the argument.} \\
\\ 
\texttt{Trial \#2: New plan: } \\
\\ 
\texttt{I failed to identify the validity of the fourth argument because I didn't properly analyze the relationship between the premises and the conclusion. Upon re{-}examining the argument, I realized that the conclusion "Being an ancestor of Janice is necessary for being a close friend of Francine" does not logically follow from the premises. However, I also noticed that I made a mistake in my analysis of the argument. The premises state that every close friend of Francine is a cousin of Samantha or an ancestor of Janice or a close friend of Shanika, and that being a close friend of Shanika is necessary for being a close friend of Francine. This implies that being a close friend of Francine is equivalent to being a close friend of Shanika. Therefore, the conclusion "Being a close friend of Francine is equivalent to being a close friend of Shanika" logically follows from the premises. However, this conclusion is not the one I was supposed to arrive at. I should have noticed that the premises do not provide any information about the relationship between being a close friend of Francine and being an ancestor of Janice. Next time, I will carefully analyze the premises and the conclusion to determine the validity of the argument, and I will also pay attention to the specific conclusion I am supposed to arrive at.} \\
\\ 
\texttt{Trial \#3: New plan: } \\
\\ 
\texttt{I failed to identify the validity of the fourth argument because I didn't properly analyze the relationship between the premises and the conclusion. Upon re{-}examining the argument, I realized that the conclusion "Being an ancestor of Janice is necessary for being a close friend of Francine" does not logically follow from the premises. However, I also noticed that I made a mistake in my analysis of the argument. The premises state that every close friend of Francine is a cousin of Samantha or an ancestor of Janice or a close friend of Shanika, and that being a close friend of Shanika is necessary for being a close friend of Francine. This implies that being a close friend of Francine is equivalent to being a close friend of Shanika. Therefore, the conclusion "Being a close friend of Francine is equivalent to being a close friend of Shanika" logically follows from the premises. However, this conclusion is not the one I was supposed to arrive at. I should have noticed that the premises do not provide any information about the relationship between being a close friend of Francine and being an ancestor of Janice. Next time, I will carefully analyze the premises and the conclusion to determine the validity of the argument, and I will also pay attention to the specific conclusion I am supposed to arrive at. To improve my analysis, I will make sure to identify the specific relationships between the premises and the conclusion, and I will also pay attention to the specific conclusion I am supposed to arrive at.} \\
\\ 
\\ 
\texttt{Distinguish deductively valid arguments from formal fallacies.} \\
\\ 
\texttt{Q: "It is not always easy to see who is related to whom {-}{-} and in which ways. The following argument pertains to this question: To begin with, Lesley is a close friend of Fernando. Moreover, being a close friend of Fernando or a schoolmate of Lowell is sufficient for being a great{-}grandfather of Leroy. It follows that Lesley is a great{-}grandfather of Leroy."} \\
\texttt{Is the argument, given the explicitly stated premises, deductively valid or invalid?} \\
\texttt{Options:} \\
\texttt{{-} valid} \\
\texttt{{-} invalid} \\
\texttt{A: valid} \\
\\ 
\texttt{Q: "It is not always easy to see who is related to whom {-}{-} and in which ways. The following argument pertains to this question: Whoever is not a great{-}grandfather of Clyde is a stepbrother of Brian. Being an ancestor of Dana is sufficient for not being a great{-}grandfather of Clyde. We may conclude: Everyone who is an ancestor of Dana is a stepbrother of Brian, too."} \\
\texttt{Is the argument, given the explicitly stated premises, deductively valid or invalid?} \\
\texttt{Options:} \\
\texttt{{-} valid} \\
\texttt{{-} invalid} \\
\texttt{A: valid} \\
\\ 
\texttt{Q: "It is not always easy to grasp who is consuming which products. The following argument pertains to this question: Every infrequent user of Paul Mitchell shampoo is either a rare consumer of Nioxin shampoo or a loyal buyer of Caress soap, or both. No regular consumer of Lush soap is a rare consumer of Nioxin shampoo and, in the same time, a loyal buyer of Caress soap. It follows that whoever is an infrequent user of Paul Mitchell shampoo is not a regular consumer of Lush soap."} \\
\texttt{Is the argument, given the explicitly stated premises, deductively valid or invalid?} \\
\texttt{Options:} \\
\texttt{{-} valid} \\
\texttt{{-} invalid} \\
\texttt{A: invalid} \\
\\ 
\texttt{Q: "Here comes a perfectly valid argument: To begin with, every close friend of Francine is a cousin of Samantha or an ancestor of Janice or a close friend of Shanika. Moreover, every close friend of Francine is a cousin of Samantha. Finally, being a close friend of Shanika is necessary for being a close friend of Francine. We may conclude: Being an ancestor of Janice is necessary for being a close friend of Francine."} \\
\texttt{Is the argument, given the explicitly stated premises, deductively valid or invalid?} \\
\texttt{Options:} \\
\texttt{{-} valid } \\
\texttt{{-} invalid} \\
\texttt{A:} \\
            }
        }
    }
     \\ \\

    \theutterance \stepcounter{utterance}  
    & \multicolumn{4}{p{0.6\linewidth}}{
        \cellcolor[rgb]{0.561,0.737,0.561}{
            \makecell[{{p{\linewidth}}}]{
                \texttt{\tiny{[P1$\rangle$GM]}}
                \texttt{Based on the provided information, I will analyze each argument to determine whether it is deductively valid or invalid.} \\
\\ 
\texttt{1. "It is not always easy to see who is related to whom {-}{-} and in which ways. The following argument pertains to this question: To begin with, Lesley is a close friend of Fernando. Moreover, being a close friend of Fernando or a schoolmate of Lowell is sufficient for being a great{-}grandfather of Leroy. It follows that Lesley is a great{-}grandfather of Leroy."} \\
\\ 
\texttt{This argument is deductively valid. The premises state that being a close friend of Fernando or a schoolmate of Lowell is sufficient for being a great{-}grandfather of Leroy. Since Lesley is a close friend of Fernando, it logically follows that Lesley is a great{-}grandfather of Leroy.} \\
\\ 
\texttt{2. "It is not always easy to see who is related to whom {-}{-} and in which ways. The following argument pertains to this question: Whoever is not a great{-}grandfather of Clyde is a stepbrother of Brian. Being an ancestor of Dana is sufficient for not being a great{-}grandfather of Clyde. We may conclude: Everyone who is an ancestor of Dana is a stepbrother of Brian, too."} \\
\\ 
\texttt{This argument is deductively valid. The premises state that being an ancestor of Dana is sufficient for not being a great{-}grandfather of Clyde. Since being an ancestor of Dana implies not being a great{-}grandfather of Clyde, it logically follows that everyone who is an ancestor of Dana is a stepbrother of Brian.} \\
\\ 
\texttt{3. "It is not always easy to grasp who is consuming which products. The following argument pertains to this question: Every infrequent user of Paul Mitchell shampoo is either a rare consumer of Nioxin shampoo or a loyal buyer of Caress soap, or both. No regular consumer of Lush soap is a rare consumer of Nioxin shampoo and, in the same time, a loyal buyer of Caress soap. It follows that whoever is an infrequent user of Paul Mitchell shampoo is not a regular consumer of Lush soap."} \\
\\ 
\texttt{This argument is deductively invalid. The premises do not provide any information about the relationship between being an infrequent user of Paul Mitchell shampoo and being a regular consumer of Lush soap. The conclusion does not logically follow from the premises.} \\
\\ 
\texttt{4. "Here comes a perfectly valid argument: To begin with, every close friend of Francine is a cousin of Samantha or an ancestor of Janice or a close friend of Shanika. Moreover, every close friend of Francine is a cousin of Samantha. Finally, being a close friend of Shanika is necessary for being a close friend of Francine. We may conclude: Being an ancestor of Janice is necessary for being a close friend of Francine."} \\
\\ 
\texttt{This argument is deductively invalid. The premises state that every close friend of Francine is a cousin of Samantha or an ancestor of Janice or a close friend of Shanika. However} \\
            }
        }
    }
    & & \\ \\

    \theutterance \stepcounter{utterance}  
    & & & \multicolumn{2}{p{0.3\linewidth}}{
        \cellcolor[rgb]{0.9,0.9,0.9}{
            \makecell[{{p{\linewidth}}}]{
                \texttt{\tiny{[GM$|$GM]}}
                \texttt{Target: invalid} \\
            }
        }
    }
    & & \\ \\

    \theutterance \stepcounter{utterance}  
    & & & \multicolumn{2}{p{0.3\linewidth}}{
        \cellcolor[rgb]{0.9,0.9,0.9}{
            \makecell[{{p{\linewidth}}}]{
                \texttt{\tiny{[GM$|$GM]}}
                \texttt{game\_result = LOSE} \\
            }
        }
    }
    & & \\ \\

    \theutterance \stepcounter{utterance}  
    & & \multicolumn{4}{p{0.6\linewidth}}{
        \cellcolor[rgb]{0.95,0.95,0.95}{
            \makecell[{{p{\linewidth}}}]{
                \texttt{\tiny{[P1$\leftrightarrow$P]}}
                \texttt{Post{-}episode Reflexion (not sent to GM):} \\
\texttt{New plan: } \\
\\ 
\texttt{I failed to identify the validity of the fourth argument in the past attempts because I didn't properly analyze the relationship between the premises and the conclusion. Upon re{-}examining the argument, I realized that the conclusion "Being an ancestor of Janice is necessary for being a close friend of Francine" does not logically follow from the premises. However, I also noticed that I made a mistake in my analysis of the argument. The premises state that every close friend of Francine is a cousin of Samantha or an ancestor of Janice or a close friend of Shanika, and that being a close friend of Shanika is necessary for being a close friend of Francine. This implies that being a close friend of Francine is equivalent to being a close friend of Shanika. Therefore, the conclusion "Being a close friend of Francine is equivalent to being a close friend of Shanika" logically follows from the premises. However, this conclusion is not the one I was supposed to arrive at. I should have noticed that the premises do not provide any information about the relationship between being a close friend of Francine and being an ancestor of Janice. } \\
\\ 
\texttt{To improve my analysis, I will make sure to identify the specific relationships between the premises and the conclusion, and I will also pay attention to the specific conclusion I am supposed to arrive at. I will break down the argument into smaller steps and analyze each step separately. I will also look for any implicit assumptions or unstated premises that may affect the validity of the argument. } \\
\\ 
\texttt{Specifically, I will focus on the following steps:} \\
\\ 
\texttt{1. Identify the specific relationships between the premises and the conclusion.} \\
\texttt{2. Analyze each premise separately to understand its contribution to the argument.} \\
\texttt{3. Look for any implicit assumptions or unstated premises that may affect the validity of the argument.} \\
\texttt{4. Pay attention to the specific conclusion I am supposed to arrive at and make sure it logically follows from the premises.} \\
\\ 
\texttt{By following these steps, I hope to improve my analysis and correctly identify the validity of the argument.} \\
            }
        }
    }
    & \\ \\

\end{supertabular}
}

\end{document}
